\chapter{Revisão de Literatura}

\section{Modelos Teóricos Fundamentais dos Ciclos Políticos}

A literatura sobre a interação entre política e macroeconomia tradicionalmente se divide em duas grandes vertentes de investigação: a teoria dos ciclos político-econômicos (CPE, ou \textit{political business cycles}) e a teoria dos ciclos político-orçamentários (CPO, ou \textit{political budget cycles}). Enquanto a primeira analisa a manipulação de agregados macroeconômicos reais, como o produto interno bruto (PIB), a taxa de desemprego e a inflação, a segunda foca na utilização discricionária de instrumentos de política fiscal (receitas, despesas correntes, investimentos) e na consequente recomposição de despesas públicas ao redor do calendário eleitoral.

A formulação pioneira do ciclo político oportunista de negócios foi proposta por \textcite{nordhaus1975}. O pressuposto fundamental do modelo de Nordhaus baseia-se na existência de uma curva de Phillips de curto prazo, que permite um \textit{trade-off} temporário entre inflação e desemprego. Sob a suposição de que os eleitores são dotados de expectativas adaptativas e exibem um comportamento "míope" (acabam priorizando retrospectivamente os resultados econômicos observados no período imediatamente anterior ao pleito), os governantes oportunistas detêm as condições necessárias para manipular deliberadamente a economia com fins eleitorais. Nessa estrutura, os partidos são puramente orientados para o poder. Seu único objetivo é ganhar eleições e maximizar a quantidade de votos (pluralidade) no próximo pleito, utilizando as políticas econômicas para atingir esse fim.

O funcionamento dinâmico desse ciclo ocorre de forma padronizada de acordo com o calendário das eleições. No período pós-eleitoral, logo após a eleição, o governante vitorioso adota políticas macroeconômicas austeras e contracionistas para elevar intencionalmente o desemprego a um nível relativamente alto, com o objetivo de combater a inflação herdada do período anterior. Na sequência, durante o período pré-eleitoral e com a aproximação do pleito, o incumbente reverte essa orientação, acionando estímulos monetários e fiscais expansionistas para aquecer a atividade econômica e reduzir o desemprego na véspera da eleição até o ponto puramente míope. Essa expansão artificial gera uma sensação temporária de prosperidade que atrai o voto dos eleitores, mas os custos inflacionários decorrentes dessa manipulação oportunista são postergados para o período posterior às eleições, reiniciando o padrão cíclico. Esse comportamento clássico oportunista gera flutuações e impulsos sistemáticos sobre a taxa de crescimento do PIB, do desemprego e da expansão monetária.

Para validar sua teoria, Nordhaus analisou o comportamento de variáveis econômicas em relação às datas de eleições em quatro democracias capitalistas (EUA, Reino Unido, Suécia e Nova Zelândia). Para os EUA, Reino Unido e Suécia, o autor avalia as flutuações nas taxas de desemprego. No entanto, para a Nova Zelândia, por se tratar de uma economia aberta com desemprego insignificante no período, a variável de controle analisada é a taxa de crescimento das importações. O autor encontra evidências estatisticamente significativas do ciclo político de negócios nos EUA e na Nova Zelândia (onde as importações eram restringidas no início do mandato e liberadas às vésperas do pleito para agradar o eleitor). O Reino Unido também apresentou fortes evidências do ciclo, embora o autor ressalte que crises no balanço de pagamentos restringiram a política macroeconômica britânica em momentos específicos. Por outro lado, os dados referentes à Suécia se mostraram estatisticamente fracos e inconclusivos, com as oscilações podendo ser perfeitamente atribuídas ao acaso.

Como um contraponto à perspectiva estritamente oportunista, a dimensão ideológica foi introduzida no estudo das flutuações macroeconômicas, estabelecendo as bases dos chamados Ciclos Políticos Partidários. Ao contrário da premissa de que os governantes agem de maneira idêntica na busca pela reeleição, o modelo de \textcite{hibbs1977} postula que os partidos políticos são orientados por preferências ideológicas distintas e consistentes, buscando atender aos interesses econômicos específicos de suas respectivas bases eleitorais. A conduta macroeconômica, sob essa ótica, baseia-se em escolhas divergentes no \textit{trade-off} estável entre inflação e desemprego: enquanto os partidos de esquerda priorizam o combate ao desemprego, a expansão econômica e políticas de distribuição de renda, mostrando-se dispostos a tolerar taxas de inflação mais elevadas em benefício das classes trabalhadoras, os partidos de centro e direita priorizam o controle inflacionário e a estabilidade de preços, adotando políticas de austeridade monetária e fiscal que frequentemente resultam em maiores taxas de desemprego no curto e médio prazo.

Empiricamente, utilizando séries temporais do pós-guerra (1948-1972) para os Estados Unidos e o Reino Unido, evidenciou-se que os períodos governados pela esquerda registraram taxas de desemprego sistematicamente menores, ao passo que gestões de centro e direita exibiram maior desemprego e maior rigidez no controle inflacionário. Dessa forma, as flutuações econômicas deixam de refletir meramente o calendário eleitoral e passam a derivar diretamente da alternância ideológica no poder.

Com o advento da revolução das expectativas racionais na macroeconomia, os modelos clássicos de ciclos políticos baseados na miopia sistemática dos agentes econômicos e na existência de um \textit{trade-off} explorável de longo prazo na curva de Phillips sofreram severas críticas teóricas, uma vez que agentes racionais que compreendem os incentivos governamentais seriam capazes de antecipar e neutralizar as manobras expansionistas das autoridades. Nesse contexto de transição metodológica, \textcite{alesina1987} reformulou a teoria partidária clássica ao integrar a hipótese de agentes racionais e orientados para o futuro (\textit{forward-looking}) ao modelo de diferenciação ideológica de Hibbs.

O Modelo Partidário Racional descreve a interação estratégica em um sistema bipartidário com trabalhadores e definidores de salários racionais. Como os contratos salariais nominais são firmados de forma antecipada para vigorar no período subsequente e possuem rigidez temporária, os agentes privados precisam projetar a inflação esperada para manter o salário real estável. No entanto, o resultado das eleições é intrinsecamente incerto, o que impede os agentes de saberem de antemão qual partido sairá vitorioso nas urnas.

Consequentemente, os salários nominais são ajustados com base em uma taxa de inflação média esperada que pondera a probabilidade eleitoral de cada partido concorrente: Se o partido de esquerda (mais expansionista) vence o pleito, ocorre uma expansão monetária e fiscal superior à inflação média esperada que foi refletida nos contratos vigentes, gerando uma surpresa inflacionária positiva que reduz temporariamente os salários reais e eleva o produto e o emprego acima de suas taxas naturais. Se o partido de direita (mais contracionista) vence o pleito, a inflação observada é menor do que a projetada pelos contratos nominais, gerando uma surpresa deflacionária que induz uma recessão temporária e elevação do desemprego no início do mandato.

A grande inovação teórica de \textcite{alesina1987} reside em demonstrar que as flutuações na atividade real são de caráter estritamente temporário e restritas ao início de cada administração. À medida que os contratos de trabalho antigos expiram e os agentes econômicos ajustam racionalmente suas expectativas inflacionárias à nova realidade governamental, a economia retorna à sua taxa natural de desemprego e produto, restabelecendo a neutralidade da política monetária no longo prazo.

A transição metodológica das discussões macroeconômicas de ciclos tradicionais para a análise dos Ciclos Políticos Orçamentários (CPO) sob expectativas racionais consolidou-se por meio dos trabalhos de \textcite{rogoff1988} e de \textcite{rogoff1990}. Diante da impossibilidade de enganar sistematicamente agentes racionais por meio de surpresas monetárias e fiscais repetitivas, os autores demonstraram que os ciclos eleitorais emergem a partir de assimetrias de informação temporárias entre governantes e eleitores na execução fiscal.

Nesses modelos de sinalização, os eleitores são agentes racionais, orientados para o futuro (\textit{forward-looking}), que desejam reconduzir ao poder o governante que demonstre maior competência administrativa na gestão das contas públicas. A competência do gestor - definida como a capacidade de prover um determinado nível de serviços públicos arrecadando o menor montante possível de impostos é uma informação privada conhecida contemporaneamente apenas pelo próprio incumbente. Como a sociedade observa certos resultados orçamentários, a exemplo da efetividade dos investimentos públicos de longo prazo, apenas de forma defasada, o governante detém uma vantagem informacional temporária.

Diante desse cenário, o ciclo orçamentário eleitoral surge como um equilíbrio separador impulsionado pela necessidade de sinalização. O modelo prevê que governantes competentes utilizam a distorção da política fiscal pré-eleitoral como um sinal de sua eficiência. Ao reduzir impostos e elevar gastos públicos de rápida percepção pela população, o incumbente consegue se diferenciar da oposição. O gestor incompetente, por sua vez, enfrenta custos marginais significativamente mais altos para imitar essa mesma expansão fiscal, o que inviabiliza a cópia sistemática da estratégia e resulta na revelação da verdadeira capacidade gerencial de cada candidato antes do pleito.

Dessa forma, os ciclos orçamentários eleitorais racionalizados por \textcite{rogoff1988} e \textcite{rogoff1990} não derivam da irracionalidade ou da 'miopia' passiva do eleitor. Pelo contrário, constituem um mecanismo de sinalização racionalmente aceito para mitigar o problema de seleção adversa no mercado político, ainda que ocorram à custa de distorções e realocações fiscais ineficientes no curto prazo.

Em suma, a evolução da literatura evidencia uma transição de modelos baseados na ilusão sistemática dos eleitores \parencite{nordhaus1975, hibbs1977} para abordagens fundamentadas na racionalidade econômica e em falhas de mercado, como a rigidez contratual \parencite{alesina1987} e a assimetria de informação \parencite{rogoff1990}. Juntos, esses quatro pilares teóricos fornecem o arcabouço analítico essencial para compreender como as restrições eleitorais, as divergências ideológicas e as estratégias de sinalização fiscal moldam a dinâmica macroeconômica nas democracias contemporâneas, servindo de base para as investigações empíricas apresentadas a seguir neste trabalho.

\section{Revisão de Literatura Empírica: Evidências e Desenvolvimentos Aplicados}

O avanço da literatura empírica sobre ciclos políticos buscou testar as previsões dos modelos teóricos clássicos e racionais, confrontando a hipótese de oportunismo governamental com o comportamento real do eleitorado e a alocação dos recursos públicos. Enquanto as primeiras investigações focaram em agregados macroeconômicos em nível nacional, a literatura subsequente concentrou-se na dinâmica das finanças subnacionais. O debate central gira em torno de duas frentes: se os eleitores recompensam ou punem a expansão fiscal eleitoral nas urnas e se os governantes operam por meio da expansão do déficit agregado ou se utilizam a recomposição qualitativa das despesas públicas.

A investigação sobre a eficácia eleitoral das manipulações fiscais oportunistas divide os pesquisadores entre a tese do conservadorismo fiscal e a da recompensa ao oportunismo. Na literatura internacional, \textcite{brender2005, brender2008} apresentam evidências contundentes que desafiam a visão convencional de que políticas fiscais expansionistas e a geração de déficits agregados auxiliam na reeleição dos governantes. Ao analisarem um amplo painel de países democráticos, os autores demonstram que os déficits em anos eleitorais não elevam as chances de reeleição; pelo contrário, o comportamento deficitário pode até punir o governante nas urnas. Esse fenômeno de conservadorismo fiscal faz com que os ciclos políticos orçamentários de déficits agregados fiquem restritos a democracias jovens e menos consolidadas, onde os eleitores possuem menor experiência histórica com eleições livres.

Ademais, \textcite{shi2006} explicam que a magnitude do risco moral envolvido nas manipulações fiscais é condicionada pela qualidade das instituições e da governança. O ciclo orçamentário tende a ser mais acentuado em países em desenvolvimento devido à combinação de duas variáveis fundamentais: o elevado montante de rendas decorrentes de se manter no poder e o menor percentual de eleitores informados na população, que enfrentam dificuldades para auditar o orçamento e monitorar de forma eficiente a conduta dos governantes.

No Brasil, os testes empíricos em nível estadual revelam uma dinâmica distinta daquela encontrada internacionalmente. Ao analisarem o comportamento fiscal e eleitoral dos estados brasileiros entre 1986 e 2002, \textcite{nakaguma2010} propõem uma metodologia de decomposição dos ciclos políticos para separar a parcela devida ao oportunismo fiscal da parcela decorrente da atividade de sinalização de competência. Ao contrário do conservadorismo fiscal documentado nas democracias maduras, o resultado principal indica que o eleitorado estadual brasileiro recompensa fortemente a parcela oportunista dos ciclos eleitorais, tanto nas receitas quanto nas despesas, embora a magnitude deste efeito tenha registrado uma trajetória de declínio progressivo ao longo do tempo, apresentando uma nítida tendência de queda ao comparar o período de 1990-1994 com 1998-2002.

No âmbito municipal, os trabalhos brasileiros também encontraram evidências de recompensa ao oportunismo e limitações na racionalidade perfeita do eleitorado. \textcite{sakurai2005} investigou a existência de ciclos orçamentários racionais em um painel de 572 municípios paulistas entre 1989 e 2001. Utilizando uma abordagem metodológica baseada no teste de hipóteses de ciclos racionais, o autor identificou impulsos fiscais significativos nos anos eleitorais, demonstrando que os prefeitos paulistas que conseguem a reeleição despendem, a princípio, um volume adicional de recursos em anos de pleito. Contudo, ao estimar a probabilidade de reeleição especificamente para o pleito municipal de 2000 (o primeiro a permitir a reeleição direta dos prefeitos municipais no Brasil), \textcite{sakurai2005} demonstrou que apenas os gastos previamente conhecidos e previstos pelos eleitores afetaram positivamente as chances de permanência do governante, ao passo que os choques de competência (resíduos) mostraram-se estatisticamente nulos. Os testes de igualdade de coeficientes indicaram que os eleitores não conseguiram distinguir de forma clara o sinal (competência) do ruído (oportunismo), o que rejeita a hipótese de racionalidade do eleitorado nos moldes tradicionais de sinalização.

Ampliando essa perspectiva empírica paulista, \textcite{sakurai2007} investigaram a conduta orçamentária de 572 prefeituras entre 1989 e 2001. Utilizando regressões de efeitos fixos, os autores identificaram impulsos fiscais oportunistas significativos nas despesas orçamentárias totais e despesas correntes nos pleitos de 1992 e 1996. Paralelamente, o estudo trouxe importantes evidências sobre a existência de ciclos partidários setoriais no nível local: prefeituras geridas pelo PFL apresentaram volumes de investimentos em capital público significativamente maiores do que os demais partidos, enquanto gestões do PTB, PPB/PDS e PMDB demonstraram uma preferência sistemática pela alocação de recursos em transferências correntes a terceiros. Por fim, \textcite{sakurai2007} testaram econometricamente o impacto imediato da Lei de Responsabilidade Fiscal, identificando que a introdução da nova norma contábil impôs uma severa contração fiscal média de R\$ 60,55 \textit{per capita} na despesa orçamentária agregada já nos anos de 2000 e 2001, forçando uma retração coordenada tanto nos gastos de consumo corrente quanto nos investimentos locais de capital.

Ampliando a abrangência geográfica, \textcite{sakurai2008} analisaram um painel de 2.235 municípios brasileiros entre 1988 e 2003. Os resultados corroboraram a tese de que a expansão fiscal é recompensada na realidade brasileira, comprovando que a elevação das despesas públicas municipais de fato eleva de forma estatisticamente significante as probabilidades de reeleição do prefeito ou de seu sucessor partidário. De maneira mais detalhada, os autores constataram que os canais fiscais eleitoralmente eficazes dependem da natureza do gasto: enquanto a manutenção de despesas de capital (investimentos) elevadas ao longo dos anos anteriores ao pleito aumenta as chances de reeleição, no caso das despesas correntes o ganho eleitoral decorre de desvios expansionistas oportunistas realizados especificamente no ano do pleito.

Em uma análise de abrangência nacional subsequente, \textcite{sakurai2011} exploraram as dinâmicas de ciclos oportunistas e partidários ao avaliarem um painel composto por 2.527 municípios brasileiros entre 1989 e 2005. Para mitigar potenciais vieses de endogeneidade da dinâmica fiscal, os autores aplicaram o estimador GMM em sistema de Blundell e Bond. Os resultados confirmaram a robustez do comportamento oportunista agregado no nível subnacional brasileiro, revelando um aumento sistemático nas despesas totais e correntes acompanhado de uma contração intencional na arrecadação de tributos locais, nos investimentos e no superávit fiscal durante os anos eleitorais. Sob a ótica partidária, \textcite{sakurai2011} evidenciaram que partidos de esquerda apresentaram maior superávit fiscal e um controle estatisticamente significante sobre as despesas totais em relação às coalizões de centro. No tocante às despesas correntes, contudo, o coeficiente associado à esquerda não se mostrou estatisticamente significante, o que corrobora a conclusão dos autores de que as despesas de custeio corrente municipal são correlacionadas de forma fraca com a ideologia do governante, concentrando-se o ajuste fiscal ideológico na contenção do gasto público agregado. Adicionalmente, os autores documentaram que o alinhamento político com o governo estadual funciona como um canalizador de recursos condicional: a aliança partidária entre o prefeito e o governador eleva as despesas totais e correntes especificamente no ano de eleição municipal, indicando que as transferências voluntárias estaduais são direcionadas de forma politicamente motivada às vésperas do pleito local.

A constatação de \textcite{sakurai2011} de que os investimentos municipais sofrem uma contração sistemática em anos de eleição no Brasil estabelece um contraste empírico relevante em duas dimensões. Primeiro, em âmbito regional, diverge do padrão observado no Estado de São Paulo por \textcite{sakurai2007}, que não identificaram oscilações estatisticamente significantes nos investimentos associadas ao calendário eleitoral paulista, evidenciando que amostras mais agregadas capturam dinâmicas cíclicas distintas. Segundo, em perspectiva internacional, o resultado brasileiro contrapõe-se diretamente à tese de recomposição orçamentária oportunista de \textcite{drazen2010}, cujos achados na Colômbia apontam para uma forte expansão pré-eleitoral de investimentos em infraestrutura altamente visíveis para o eleitor (\textit{voter-friendly}). No contexto brasileiro de dados nacionais, a queda nos investimentos no ano do pleito sugere que despesas de capital demandam maior tempo de planejamento e execução, estimulando os prefeitos a buscarem apelo eleitoral imediato por meio do aumento discricionário de gastos correntes de consumo.

No que concerne ao pioneirismo na investigação das flutuações macroeconômicas nacionais no país, destaca-se o estudo clássico de \textcite{fialho1997}. A autora utilizou séries temporais para o período de 1953 a 1995, combinando premissas de expectativas adaptativas do modelo clássico com os novos modelos de expectativas racionais. O estudo revelou que o período eleitoral afeta de forma positiva e estatisticamente significante a dinâmica macroeconômica brasileira, gerando impactos expansionistas na taxa anual de crescimento do PIB real (que tendeu a ser cerca de 2\% maior em anos eleitorais) e na oferta real de moeda (M1). Em contrapartida, \textcite{fialho1997} demonstrou que as taxas de desemprego e inflação não registraram variações eleitorais significantes, confirmando que a inflação brasileira no período exibia forte inércia estrutural. O estudo conclui que o ciclo político sobre o produto e a moeda constitui um fenômeno estatisticamente robusto apenas a partir da redemocratização e abertura política do país iniciada em 1975.

A constatação de que eleitores podem punir desequilíbrios agregados e a vigência de restrições fiscais direcionaram o foco das análises para as manipulações qualitativas da estrutura do orçamento. Historicamente, o modelo pioneiro de sinalização de \textcite{rogoff1990} sugeria que governantes oportunistas tendiam a sinalizar competência gerencial expandindo despesas visíveis antes das eleições. No entanto, para acomodar a existência de eleitores fiscalmente conservadores e avessos ao deficit, a literatura avançou em outra direção. A principal referência conceitual desse movimento é o trabalho de \textcite{drazen2010}, que propõe a teoria da recomposição oportunista das despesas públicas. Segundo os autores, para sinalizar que compartilham das mesmas preferências e prioridades de gastos dos cidadãos no período pós-eleitoral, sem incorrer na geração de déficits agregados, os governantes alteram a composição do orçamento: reduzem gastos não direcionados diretamente aos eleitores e expandem as despesas com bens públicos de alta visibilidade direta (\textit{voter-friendly spending}).

No plano empírico, ao analisarem as finanças públicas locais de todos os municípios da Colômbia entre os anos de 1987 e 2002, \textcite{drazen2010} forneceram um respaldo robusto para esse modelo. Os resultados econométricos evidenciaram uma nítida mudança na composição orçamentária pré-eleitoral: os governantes locais expandiram significativamente os investimentos físicos em infraestrutura de alta visibilidade direta para a população (como infraestrutura urbana, habitação, saúde, água e energia), enquanto contraíram de forma coordenada despesas menos salientes ou de menor apelo político imediato, como despesas com pessoal temporário e transferências correntes (especialmente para aposentados). Ademais, o estudo comprovou que, embora os eleitores se comportem como conservadores fiscais avessos a déficits, eles respondem favoravelmente à recomposição discricionária do orçamento, recompensando nas urnas o partido do prefeito de turno que direciona recursos para esses investimentos de alta visibilidade.

No Brasil, essa hipótese de recomposição e saliência fiscal encontrou robusto respaldo empírico em diferentes escalas da federação. \textcite{sakurai2009} analisou a alocação de despesas de mais de 5.500 municípios brasileiros divididas por funções orçamentárias entre 1990 e 2005. Os resultados evidenciaram que o ano de eleições municipais afeta diretamente a priorização funcional do orçamento municipal, gerando expansões sistemáticas e estatisticamente significativas em áreas de forte apelo social e visual direto, tais como saúde e saneamento, habitação e urbanismo, assistência e previdência, e transportes. Em contrapartida, funções orçamentárias menos salientes ou de retorno político menos imediato, como a agricultura, sofrem retração no período eleitoral, consolidando a evidência empírica de um nítido efeito composição impulsionado pelo calendário eleitoral.

Analisando municípios brasileiros em um contexto institucional pós-Lei de Responsabilidade Fiscal (LRF), \textcite{nunes2017} investigou as despesas sociais de saúde e educação, bem como os investimentos de capital, em 407 municípios gaúchos entre 2002 e 2012. A análise econométrica de painel com efeitos fixos revelou que as despesas per capita em saúde e educação não registraram elevações estatisticamente significantes em anos eleitorais gaúchos. De acordo com o autor, a rigidez orçamentária decorrente de vinculações constitucionais de receitas e os limites rígidos impostos pela LRF impedem o remanejamento discricionário dessas áreas sociais de forma oportunista no curto prazo. Entretanto, a variável ano eleitoral apresentou impacto positivo, expressivo e altamente significante sobre as despesas com investimentos de capital, confirmando que os prefeitos gaúchos recorrem à recomposição orçamentária (priorizando as obras públicas visíveis em detrimento de despesas correntes inelásticas) para sinalizar competência ao eleitorado em conformidade com as regras fiscais vigentes.

Adicionalmente, expandindo a discussão sobre como a intensidade da disputa eleitoral afeta essa dinâmica fiscal, \textcite{teixeira2021} investigaram o impacto da competição política sobre a alocação dos recursos nos municípios brasileiros entre 2004 e 2016. Os autores introduzem a discussão sobre a saliência dos gastos, argumentando que os prefeitos enfrentando cenários de maior concorrência eleitoral tendem a estruturar o orçamento de forma estratégica no ano da eleição. Os resultados empíricos mostram que, sob maior competição política, os governantes optam por reduzir proporcionalmente os gastos administrativos locais (considerados não salientes e de baixa percepção pelo eleitorado). Esse movimento permite poupar recursos em rubricas de baixa visibilidade para manter ou elevar as despesas mais perceptíveis em períodos eleitorais, corroborando a ideia de que o acirramento das eleições atua como um catalisador de comportamentos estratégicos de recomposição qualitativa dos orçamentos públicos.

\section{Mecanismos Específicos, Restrições Institucionais e Dinâmicas Espaciais no Brasil}

A investigação empírica dos Ciclos Políticos Orçamentários (CPO) no Brasil avançou substancialmente ao incorporar as especificidades do arranjo federativo do país, a existência de regras fiscais rígidas e a utilização estratégica de mecanismos contábeis discricionários. Longe de se limitar a uma expansão linear e homogênea dos déficits agregados, a manipulação eleitoral no contexto brasileiro é caracterizada por comportamentos fiscais sofisticados, nos quais prefeitos gerenciam a visibilidade e a temporalidade de seus gastos sob restrições institucionais severas.

No federalismo fiscal brasileiro, a dependência financeira dos governos locais em relação às esferas superiores introduz um componente político central na dinâmica dos ciclos fiscais. Ao investigarem o papel das transferências voluntárias da União e dos Estados para os municípios, \textcite{ferreira2007} demonstram estatisticamente que a distribuição desses recursos não é neutra. Registra-se que os municípios recebem, em média, volumes significativamente maiores de transferências não-constitucionais quando o prefeito em exercício pertence à coligação partidária que elegeu o governador do estado ou ao mesmo partido político do Presidente da República.

Ao estender o modelo clássico de sinalização de \textcite{rogoff1990} para contemplar dois níveis de governo com eleições intercaladas, os autores revelam que essas transferências politicamente motivadas alteram profundamente os incentivos eleitorais locais. Em um cenário de informação completa sobre a capacidade gerencial dos prefeitos, o alinhamento político funciona como uma fricção orçamentária capaz de distorcer o mecanismo de seleção de governantes. Diante do expressivo acréscimo de recursos trazido pelo apoio político do Estado, os eleitores podem, de forma estritamente racional, preferir reeleger um prefeito considerado incompetente, mas alinhado ao governador, em detrimento de um candidato opositor com maior competência gerencial esperada, mas que governaria sem o fluxo de transferências voluntárias. De maneira inversa, o eleitorado pode deixar de reconduzir um governante local competente caso o oponente seja o aliado preferencial do chefe do Executivo estadual.

Sob a premissa de assimetria informacional, na qual os eleitores não observam diretamente a competência gerencial corrente do gestor antes do pleito, a interação se configura como um clássico jogo de sinalização. Os autores demonstram que os equilíbrios agregadores são inconsistentes sob critérios de refinamento teórico, restando apenas os equilíbrios separadores como soluções válidas. A dinâmica desses equilíbrios separadores passa a ser inteiramente governada pelo montante de transferências voluntárias, dividindo-se em dois regimes distintos.

No primeiro regime, caracterizado por transferências voluntárias pouco expressivas, os eleitores continuam priorizando a seleção do candidato de maior competência administrativa. Nesse contexto, o ciclo político-orçamentário emerge como um mecanismo de sinalização custosa: prefeitos competentes são forçados a distorcer as contas públicas em ano eleitoral, reduzindo impostos ou inflando gastos visíveis contemporâneos, para comprovar sua eficiência gerencial e garantir a reeleição. Contudo, a presença das transferências voluntárias amplia a intensidade de distorção do ciclo fiscal. Ademais, sob certas condições, essas transferências podem criar artificialmente um ciclo político-orçamentário que não ocorreria na ausência delas. Ao flexibilizar a restrição orçamentária dos prefeitos aliados, o repasse discricionário reduz a disparidade prática na provisão de bens públicos entre governantes eficientes e ineficientes. Consequentemente, para conseguir se diferenciar de um gestor incompetente (cuja capacidade financeira aparente foi inflada pelas verbas extras do governador), o prefeito competente é empurrado a incorrer em distorções fiscais ainda mais agressivas de sinalização.

Por outro lado, instala-se um segundo regime quando o volume das transferências politicamente motivadas é altamente expressivo. Nesse cenário, as transferências são amplas o suficiente para compensar qualquer ineficiência administrativa, fazendo com que os eleitores prefiram o alinhamento partidário a qualquer critério de competência gerencial. Em equilíbrio, o ciclo político-orçamentário deixa de existir, o que elimina o risco moral fiscal, uma vez que nenhum governante passa a ter incentivos para distorcer as contas com finalidades de sinalização de sua competência. Em contrapartida, esse processo gera um severo quadro de seleção adversa, no qual o governador do estado passa a definir indiretamente os pleitos municipais por meio do repasse discricionário de verbas. Isso resulta na permanência sistemática de políticos gerencialmente incompetentes no poder local, o que prejudica gravemente o desenvolvimento institucional dos municípios no longo prazo.

A conduta fiscal oportunista em ano de eleição não ocorre de maneira isolada em cada jurisdição municipal, sendo influenciada pelas escolhas orçamentárias adotadas pelos municípios limítrofes. A inserção da dimensão espacial na literatura de finanças públicas baseia-se na premissa de que os governos locais sofrem forte influência de seus vizinhos ao definirem suas políticas de gastos. Essa interação horizontal ocorre por meio de dois mecanismos centrais: o efeito de transbordamento (spillover) e a competição por comparação (yardstick competition). No spillover, as decisões de despesa de uma municipalidade geram externalidades físicas diretas para as regiões vizinhas, podendo funcionar de forma complementar ou substituta (como no caso de investimentos compartilhados em saúde ou infraestrutura viária lindeira). Já a yardstick competition configura-se como uma resposta político-eleitoral sob assimetria de informação: os eleitores utilizam o desempenho fiscal dos municípios vizinhos como um "padrão de medida" para avaliar a competência de seu próprio governante. Sabendo que estão sob constante comparação, os prefeitos tendem a imitar o padrão de gastos da vizinhança para sinalizar eficiência administrativa e maximizar suas chances de reeleição.

Para testar empiricamente essas dinâmicas no cenário nacional, \textcite{videira2011} investigaram as despesas de investimentos, saúde e educação dos municípios brasileiros no período de 1997 a 2008. Para lidar com a endogeneidade decorrente da simultaneidade das decisões fiscais entre os entes vizinhos, os autores adotaram uma modelagem de dados em painel com efeitos fixos e variáveis instrumentais pautadas nas características demográficas da vizinhança. Os resultados econométricos demonstram que a inclusão do componente espacial altera profundamente a interpretação tradicional dos ciclos políticos oportunistas. Em modelos que desconsideram o espaço, o ano de eleição aparece associado a uma elevação geral de gastos. No entanto, ao controlar as estimações pela média de despesas da microrregião, o efeito isolado do ano de pleito perde significância ou inverte seu sinal para negativo, evidenciando que a expansão de despesas em períodos eleitorais ocorre de forma relativa e condicional às ações dos vizinhos.

Ao analisar as categorias individualmente, observa-se que os gastos com investimentos e educação apresentam interações espaciais positivas que são significativamente potencializadas nos anos de eleição, o que corrobora a hipótese de que os prefeitos reagem de forma estratégica e competitiva à conduta fiscal da vizinhança imediata para atrair votos. Por outro lado, a despesa com saúde apresenta um forte efeito de vizinhança direto, mas essa dinâmica de interação não se altera sob a motivação do calendário eleitoral, uma vez que a variável de interação política não apresentou significância estatística para essa função. Embora o estudo aponte para a limitação metodológica de não conseguir isolar perfeitamente qual dos dois mecanismos prevalece (se o de transbordamento físico ou o de comparação política), fica nítida a importância das interações espaciais na determinação da conduta fiscal dos prefeitos, indicando que o clássico oportunismo eleitoral nos municípios do país deve ser reinterpretado sob a ótica da interdependência regional.

À medida que as restrições institucionais de controle do endividamento público se consolidaram, os incentivos oportunistas migraram para manobras qualitativas de natureza contábil. \textcite{vicente2012} investigam de que forma as demonstrações contábeis revelam a eficácia de práticas de ilusão fiscal na probabilidade de reeleição de prefeitos ou na recondução de seus partidos de coalizão.

O comportamento contábil oportunista manifesta-se de forma saliente no gerenciamento da conta de Restos a Pagar Não Processados. Essa prática consiste em inscrever no encerramento do exercício despesas que percorreram apenas a fase de empenho (reserva orçamentária), sem que o bem tenha sido entregue ou o serviço liquidado. Essa manobra contábil distorce as demonstrações orçamentárias ao registrar formalmente despesas de investimento em infraestrutura que, na realidade física do município, não ocorreram, gerando uma ilusão de realização que atrai o eleitorado menos informado.

Empiricamente, as estimações logísticas de \textcite{vicente2012} revelam que o eleitorado municipal pune a austeridade de curto prazo: a elevação do Resultado Orçamentário Corrente (ROC) em anos eleitorais (decorrente de maior arrecadação ou de contenção de gastos com pessoal e manutenção administrativa) reduz de forma significante as chances de reeleição. Em contrapartida, os prefeitos utilizam estrategicamente a folga financeira: a obtenção de resultados financeiros positivos em anos pré-eleitorais e a geração de caixa pós-eleitoral elevam a probabilidade de recondução. O superávit financeiro acumulado no ano pré-eleitoral é intencionalmente convertido em uma maior inscrição de restos a pagar não processados em ano de pleito, permitindo inflar a provisão de serviços visíveis nas vésperas da votação.

Aprofundando os incentivos derivados das regras contábeis, \textcite{almeida2023} analisam o impacto da limitação de mandatos (primeiro mandato, elegível à reeleição, versus segundo mandato, impedido de concorrer) sobre a inscrição de restos a pagar nos municípios do país. A análise conceitua que a inscrição de Restos a Pagar Processados, que representam despesas empenhadas e efetivamente liquidadas (serviços prestados ou bens entregues), mas cujos pagamentos financeiros são postergados para o mandato seguinte, constitui um artifício de sinalização fiscal eficiente e de baixo custo, pois permite que prefeitos em primeiro mandato expandam fisicamente as obras públicas sem violar os limites de endividamento da LRF.

Os testes empíricos baseados em dados em painel e regressão em descontinuidade (RDD) trazem resultados robustos sobre essa conduta estratégica, demonstrando que prefeitos em primeiro mandato mantêm estoques e fluxos de restos a pagar processados per capita sistematicamente maiores do que prefeitos em segundo mandato. Esse comportamento é evidenciado pelo incremento no estoque médio de restos a pagar processados, que varia de 10,51\% a 11,56\% nas estimações de painel e é confirmado pelas estimativas cúbicas de RDD, indicando que prefeitos elegíveis à reeleição exibem um estoque de restos a pagar R\$ 7,90 per capita superior. Contudo, em anos eleitorais, observa-se uma retração média no estoque e no fluxo de inscrição de restos a pagar por parte de prefeitos em primeiro mandato, representada por uma redução de R\$ 6,79 per capita no fluxo de efeitos fixos e de R\$ 6,78 per capita em RDD. Esse decréscimo pré-eleitoral reflete a necessidade contábil de liquidar passivos e adequar as contas municipais às restrições do artigo 42 da LRF, que veda a contração de despesas sem disponibilidade de caixa nos últimos dois quadrimestres de mandato.

A premissa de que a inscrição contábil atende a incentivos oportunistas é validada de forma empírica através de modelos que comprovam a sua efetividade eleitoral. O modelo logit indica uma relação estatisticamente significante e positiva entre o fluxo acumulado de restos a pagar ao longo do mandato e a probabilidade de o prefeito registrar sua candidatura para a reeleição. Além disso, os prefeitos em primeiro mandato que de fato concorrem à reeleição expandem significativamente os investimentos públicos físicos em anos de eleição, registrando um incremento de R\$ 11,41 per capita em despesas de capital. Esses resultados em conjunto sustentam a hipótese de que a gestão fiscal subnacional é fortemente guiada pelos interesses de manutenção no poder e pela dinâmica dos ciclos políticos locais.

Os incentivos para a manipulação qualitativa das despesas públicas também variam de acordo com o grau de competição política enfrentado pelo governante e pelo tamanho do eleitorado do município gerido. \textcite{teixeira2021} investigam este comportamento utilizando uma metodologia de regressão em descontinuidade (sharp RDD) baseada no limiar de 200 mil eleitores registrados. No sistema eleitoral brasileiro, esse limiar determina de forma exógena a estrutura da eleição municipal: cidades abaixo deste ponto decidem o pleito em turno único (sistema majoritário plural), enquanto cidades acima adotam o sistema de dois turnos, caso nenhum candidato alcance a maioria absoluta dos votos no primeiro turno.

O modelo de regressão em descontinuidade (sharp RDD) proposto por \textcite{teixeira2021} demonstra que o formato da votação altera sensivelmente tanto a dinâmica de competição eleitoral quanto a resposta fiscal dos incumbentes municipais. No campo político, embora vigore o incentivo ao voto estratégico em municípios com turno único, os autores mostram que, para o período de 2004 a 2016, essa migração de votos para os candidatos favoritos é mais sutil do que o reportado por estudos anteriores. Especificamente, após a inclusão de controles, o efeito do voto estratégico perde significância estatística na participação do terceiro colocado em diante (lower3), mantendo-se robusto apenas a partir do quarto colocado (lower4). Paralelamente, em resposta a essa pressão competitiva, os prefeitos gerenciam de forma estratégica a composição de seus orçamentos em anos eleitorais. Essa gestão resulta em uma redução significativa na participação das despesas administrativas (consideradas operacionais e pouco salientes para o eleitorado), que registram uma retração de 2,1 a 2,2 pontos percentuais nos modelos sem controles e de 1,6 a 1,8 ponto percentual após a inserção de variáveis de controle. Esse recuo, contudo, não se traduz em aumentos estatisticamente significativos nas grandes funções sociais tradicionais, como saúde e educação. Em vez disso, as evidências secundárias e menos precisas dos testes de robustez sugerem que os recursos poupados podem ter sido redirecionados para despesas de menor representatividade, mas com forte potencial de sinalização política e apelo visual direto, tais como assistência social e transportes.

Por fim, a literatura nacional documenta uma forte heterogeneidade do comportamento fiscal oportunista conforme o tamanho financeiro e administrativo dos entes federativos. \textcite{gralak2023} examinam essa dinâmica ao investigarem a ocorrência de ciclos políticos orçamentários (CPO) nos gastos com investimentos públicos em 35 grandes municípios brasileiros (com população superior a 500 mil habitantes, excluindo-se os municípios de São Paulo e do Rio de Janeiro) no período de 2005 a 2016. A modelagem de dados em painel por efeitos aleatórios revelou que a variável relacionada às eleições municipais apresentou significância estatística apenas marginal ao nível de 10\%. Em razão desse resultado, os autores concluem categoricamente pela ausência de comportamento eleitoral oportunista cíclico nos gastos locais de capital.

Essa constatação contrariou as premissas iniciais do próprio estudo, que supunha que os grandes polos eleitorais poderiam incentivar distorções fiscais devido ao acirramento partidário pelo controle das prefeituras. No entanto, o estudo demonstra que a robustez orçamentária dessas cidades altera esse comportamento: as receitas tributárias próprias respondem por uma média de 24,72\% da arrecadação total nesses entes, conferindo-lhes expressiva autonomia e resultando em uma gestão fiscal mais preocupada com o equilíbrio orçamentário, contrariando o pressuposto de \textcite{rogoff1990}. Em vez de oscilações eleitoreiras de curto prazo, os verdadeiros impulsionadores dos investimentos municipais identificificados na amostra foram o volume de receitas tributárias, as transferências de convênios federais, a capacidade de endividamento de longo prazo (exigível a longo prazo) e o alinhamento partidário do prefeito com o governador do estado.

Essa dinâmica contrasta diretamente com a realidade fiscal pós-LRF dos municípios de pequeno e médio porte, como os analisados por \textcite{nunes2017} no Rio Grande do Sul. Ao modelar os gastos per capita com saúde, educação e investimentos em 407 municípios gaúchos entre 2002 e 2012, \textcite{nunes2017} demonstra de que forma as restrições institucionais alteraram as escolhas alocativas locais. O estimador de efeitos fixos revela que a variável ano de eleição municipal apresentou impacto negativo e altamente significante (-R\$ 49,02 per capita, significante a 1\%) sobre os gastos per capita com educação, mostrando-se estatisticamente nulo para as despesas com saúde. Como as normas de controle fiscal e as vinculações constitucionais de receitas tributárias mínimas criam uma inelasticidade orçamentária nessas funções sociais, os prefeitos gaúchos recorrem a uma recomposição qualitativa das despesas: o ano eleitoral exerce um impacto positivo e estatisticamente significante de R\$ 41,37 per capita sobre as despesas com investimentos públicos de capital, que possuem maior imediatismo e visibilidade para a população (como o planejamento e a execução de obras). Assim, a rigidez institucional decorrente da LRF e da legislação eleitoral não extinguiu o ciclo eleitoral orçamentário no âmbito municipal, mas induziu os governantes locais a realizarem uma recomposição qualitativa das despesas reduzindo temporariamente os gastos com custeio social para financiar a expansão de investimentos físicos em período pré-eleitoral. Adicionalmente, \textcite{nunes2017} identifica a existência de ciclos partidários, observando que prefeitos de esquerda despendem valores significativamente maiores em educação e saúde se comparados aos de centro.